\chapter{Esercitazione E1: reti, stratificazione e livello applicativo}
\titolobreve{Esercitazione E1}

Questo capitolo e i tre successivi raccolgono i \textbf{quattro pacchetti di esercizi} del corso (E1--E4), nell'ordine delle slide:
l'esercizio $n$ di questo capitolo è l'\textbf{E1.$n$}. Nelle slide gli esercizi non hanno soluzioni (``un risultato che non potete
confrontare con un numero è un risultato che dovete argomentare''); qui per ciascuno c'è uno svolgimento, da leggere \emph{dopo} averci
provato. Per gli esercizi di laboratorio lo svolgimento dice che cosa ci si deve aspettare: i numeri esatti dipendono dalla vostra rete.

\info{Strumenti}{
    Wireshark installato e configurato per catturare come utente normale (si avvia la cattura \emph{prima} del comando), e da terminale
    \texttt{nc}, \texttt{curl}, \texttt{dig} e \texttt{ftp}. Su Debian, Ubuntu e derivate:
    \texttt{sudo apt install netcat-openbsd curl bind9-dnsutils tnftp}. Alcune reti bloccano FTP o la porta 80 in uscita: in quel caso si
    prova da casa. I server cambiano: se un nome non risponde più, se ne sceglie un altro che faccia la stessa cosa.
}

% ══════════════════════════════════════════════════════════════
\section{Reti e commutazione}
% ══════════════════════════════════════════════════════════════

\begin{esercizio}{Host ed end system}
Qual è la differenza tra un host e un end system? Elencare diversi tipi di end system. Un server Web è un end system?

\svolgimento
Nessuna: i due termini sono sinonimi, e indicano i dispositivi ai bordi della rete che eseguono le applicazioni. Sono end system PC,
portatili, smartphone, server, sensori, televisori, console, automobili connesse. Un server Web è un end system (un host), come il
browser che lo interroga.
\end{esercizio}

\begin{esercizio}{Circuiti, TDM e FDM}
Quale vantaggio ha una rete a commutazione di circuito rispetto a una a commutazione di pacchetto? In una rete a circuito, quali
vantaggi ha il TDM rispetto all'FDM?

\svolgimento
Il circuito \textbf{garantisce} la banda per tutta la durata della chiamata: niente code, ritardo costante, nessuna perdita per
congestione, ottimo per traffico continuo come la voce. Il TDM rispetto all'FDM: ogni circuito usa tutta la banda per il proprio slot,
non serve dividere lo spettro con filtri e bande di guardia, è naturalmente digitale e si adatta meglio a canali di velocità diverse
(più slot per lo stesso circuito).
\end{esercizio}

\begin{esercizio}{Peering e IXP}
Perché due ISP allo stesso livello della gerarchia fanno spesso peering tra loro? Come guadagna un IXP?

\svolgimento
Scambiandosi direttamente il traffico tra i rispettivi clienti, senza passare dal provider di livello superiore, \textbf{non pagano}
il transito per quel traffico (il peering di solito è gratuito) e accorciano i percorsi. Un IXP è un punto in cui molti ISP si
incontrano e fanno peering: guadagna facendo pagare agli ISP la connessione alla propria infrastruttura di commutazione (porte, spazio,
corrente), non il traffico.
\end{esercizio}

\begin{esercizio}{Topologie, contate}
Otto host devono essere collegati tra loro.
\begin{enumerate}[label=(\alph*)]
    \item Per un bus, un anello, una stella attorno a uno switch e una maglia completa, quanti collegamenti servono? Per il bus si contano
          la dorsale e un cavo di discesa per host.
    \item Per ciascuna topologia, trovare un singolo guasto (un collegamento o un dispositivo) che isoli almeno un host dagli altri.
          Quanti host isola? C'è una topologia in cui nessun guasto di un singolo collegamento isola qualcuno?
    \item Con $n$ host, come cresce il numero di collegamenti nei quattro casi?
\end{enumerate}

\svolgimento
\begin{enumerate}[label=(\alph*)]
    \item Bus: 1 dorsale $+$ 8 discese $=$ 9. Anello: 8. Stella: 8. Maglia completa: $8 \cdot 7/2 = 28$.
    \item Bus: un cavo di discesa isola il suo host; un taglio della dorsale divide gli host in due gruppi (quanti, dipende dal punto).
          Anello: un collegamento rotto lascia una linea, e nessuno è isolato; ne servono due. Stella: un cavo isola il suo host, lo switch
          guasto isola tutti e otto. Maglia: nessun singolo collegamento isola nessuno. Anello e maglia sono quindi le topologie in cui nessun
          guasto di un singolo collegamento isola qualcuno (il guasto di un host isola, ovviamente, solo lui).
    \item Bus $n + 1$, anello $n$, stella $n$: lineare. Maglia $n(n-1)/2$: quadratica.
\end{enumerate}
\end{esercizio}

\begin{esercizio}{La legge di Metcalfe}
La legge di Metcalfe dice che il valore di una rete cresce con il quadrato del numero dei suoi utenti. Se ciascuno di $n$ utenti invia un
messaggio a ciascuno degli altri, quanti messaggi vengono inviati? La risposta sostiene la legge? Confrontarla con i collegamenti della
maglia completa dell'esercizio precedente.

\svolgimento
$n(n-1)$ messaggi: crescono come $n^2$, ed è proprio l'intuizione della legge (ogni nuovo utente può parlare con tutti quelli già
presenti). È esattamente il doppio dei collegamenti della maglia completa, $n(n-1)/2$: un collegamento serve due messaggi, uno per verso.
\end{esercizio}

\begin{esercizio}{Multiplexing statistico}
Un collegamento da 10 Mb/s è condiviso da utenti che, quando trasmettono, richiedono 200 kb/s ciascuno, ma trasmettono solo il 10\% del
tempo.
\begin{enumerate}[label=(\alph*)]
    \item Con la commutazione di circuito, quanti utenti si possono servire?
    \item Con la commutazione di pacchetto, qual è la probabilità che un dato utente stia trasmettendo?
    \item Con 120 utenti, qual è la probabilità che esattamente $n$ trasmettano insieme? E che siano 51 o più?
\end{enumerate}

\svolgimento
\begin{enumerate}[label=(\alph*)]
    \item $10 \text{ Mb/s} / 200 \text{ kb/s} = 50$ utenti.
    \item $p = 0{,}1$.
    \item Distribuzione binomiale: $P(n) = \binom{120}{n}\, 0{,}1^n\, 0{,}9^{120-n}$. La probabilità che ne trasmettano 51 o più (e quindi che
          il collegamento sia sovraccarico) è $\sum_{n=51}^{120} P(n) \approx 2 \cdot 10^{-20}$: in pratica mai. Con i pacchetti si servono
          120 utenti dove il circuito ne serve 50.
\end{enumerate}
\end{esercizio}

\begin{esercizio}{Direzioni, dimensioni e velocità}
Sei situazioni: (1) una tastiera senza fili invia tasti a un PC; (2) due persone parlano con i walkie-talkie; (3) una telefonata su una
linea in fibra; (4) la rete che unisce gli edifici di un'università sparsi in una città; (5) un portatile collegato allo switch
dell'ufficio; (6) due data center su continenti diversi.
\begin{enumerate}[label=(\alph*)]
    \item Per (1), (2) e (3): simplex, half-duplex o full-duplex?
    \item Per (1), (4), (5) e (6): PAN, LAN, MAN o WAN?
    \item Un download mostra 40 Mbit/s su una linea venduta come 100 Mbit/s. Quale dei due numeri è più vicino alla banda del percorso e
          quale è il throughput? Dare due ragioni per cui differiscono.
\end{enumerate}

\svolgimento
\begin{enumerate}[label=(\alph*)]
    \item (1) simplex; (2) half-duplex (si parla a turno); (3) full-duplex.
    \item (1) PAN; (4) MAN; (5) LAN; (6) WAN.
    \item 100 Mbit/s è la banda nominale del collegamento d'accesso; 40 Mbit/s è il throughput misurato. Differiscono perché il collo di
          bottiglia può stare altrove sul percorso (il server, un collegamento congestionato, il Wi-Fi di casa), perché header e
          ritrasmissioni consumano banda, perché il controllo di congestione di TCP non tiene sempre il tubo pieno, e perché il
          collegamento è condiviso con altri.
\end{enumerate}
\end{esercizio}

% ══════════════════════════════════════════════════════════════
\section{Livelli e incapsulamento}
% ══════════════════════════════════════════════════════════════

\begin{esercizio}{Chi elabora quali livelli}
Quali livelli della pila Internet elabora un router? Uno switch di collegamento? Un host?

\svolgimento
Router: fisico, collegamento e rete. Switch: fisico e collegamento. Host: tutti e cinque.
\end{esercizio}

\begin{esercizio}{Compiti dei livelli}
Elencare cinque compiti che un livello può svolgere. Lo stesso compito può essere svolto da due o più livelli? Dare un esempio dal corso.

\svolgimento
Rilevazione degli errori, controllo di flusso, segmentazione e riassemblaggio, multiplexing, indirizzamento, instaurazione della
connessione, consegna affidabile, instradamento. Sì: la rilevazione degli errori sta nel CRC del collegamento, nel checksum dell'header
IP e nel checksum di TCP/UDP; la consegna affidabile sia nel Wi-Fi (riscontri sul singolo hop) sia in TCP (da estremo a estremo).
\end{esercizio}

\begin{esercizio}{Incapsulamento a mano}
Un browser invia una richiesta HTTP di 500 byte su TCP, IPv4 ed Ethernet: header TCP di 20 byte, header IP di 20 byte, header Ethernet di
14 byte più 4 byte di controllo in fondo al frame.
\begin{enumerate}[label=(\alph*)]
    \item Disegnare il frame come scatole annidate. Quale livello aggiunge qualcosa anche in fondo?
    \item Quanti byte invia l'adattatore, senza contare il preambolo? Che frazione è la richiesta HTTP?
    \item Ripetere (b) per un messaggio applicativo di 10 byte. Che cosa dice il risultato sull'invio di molti messaggi piccoli?
    \item Un router sul percorso riceve il frame. Quali header legge, e quali riscrive prima di inoltrarlo?
\end{enumerate}

\svolgimento
\begin{enumerate}[label=(\alph*)]
    \item \texttt{[Eth 14 [IP 20 [TCP 20 [HTTP 500]]] CRC 4]}: il collegamento aggiunge anche un trailer, il CRC.
    \item $14 + 20 + 20 + 500 + 4 = 558$ byte; la richiesta è $500/558 \approx 89{,}6\%$.
    \item $14 + 20 + 20 + 10 + 4 = 68$ byte, di cui solo il $14{,}7\%$ è utile: con molti messaggi piccoli la maggior parte della banda se
          ne va in header. Conviene raggrupparli.
    \item Legge l'header Ethernet (il MAC di destinazione è il suo) e l'header IP (l'indirizzo di destinazione per la tabella). Costruisce
          un header Ethernet nuovo (MAC sorgente e destinazione, CRC) e nell'header IP riscrive TTL e checksum. TCP e HTTP non li tocca.
\end{enumerate}
\end{esercizio}

% ══════════════════════════════════════════════════════════════
\section{Applicazioni}
% ══════════════════════════════════════════════════════════════

\begin{esercizio}{Client e server nel P2P}
Per un'applicazione P2P di condivisione di file, siete d'accordo che ``non c'è nozione di lato client e lato server in una sessione di
comunicazione''? Perché?

\svolgimento
No. In ogni sessione un peer chiede (fa da client) e l'altro serve (fa da server): ciò che il P2P elimina è il ruolo \emph{fisso}, non
la distinzione in una sessione. Lo stesso peer è client in una sessione e server in un'altra, anche contemporaneamente.
\end{esercizio}

\begin{esercizio}{I servizi del trasporto}
Un protocollo di trasporto può offrire quattro tipi di servizio: consegna affidabile, throughput, temporizzazione e sicurezza. Per
ciascuno, dire se lo offre TCP, UDP, entrambi o nessuno.

\svolgimento
Consegna affidabile: solo TCP. Throughput garantito: nessuno dei due. Temporizzazione: nessuno. Sicurezza: nessuno dei due da solo; la
offre TLS sopra TCP.
\end{esercizio}

\begin{esercizio}{Progettare un protocollo applicativo}
Progettare il protocollo tra un bancomat e il computer centrale della banca: deve permettere di verificare carta e PIN, chiedere il saldo
e prelevare, compreso il caso frequente in cui il conto non basta.
\begin{enumerate}[label=(\alph*)]
    \item Elencare i tipi di messaggio, la loro sintassi, il significato di ogni campo e le regole su chi invia che cosa e quando.
    \item Disegnare lo scambio per un prelievo riuscito come diagramma temporale.
    \item Che cosa assume il protocollo sul trasporto sottostante? Che cosa succede se la richiesta del bancomat arriva ma la risposta si
          perde: il cliente riceve i soldi due volte, mai, o una volta?
\end{enumerate}

\svolgimento
\begin{enumerate}[label=(\alph*)]
    \item Messaggi di testo, una riga ciascuno: \texttt{AUTH <carta> <pin>} $\rightarrow$ \texttt{OK} oppure \texttt{ERR pin};
          \texttt{BALANCE} $\rightarrow$ \texttt{BAL <importo>}; \texttt{WITHDRAW <id> <importo>} $\rightarrow$ \texttt{GRANTED <id> <nuovo
          saldo>} oppure \texttt{DENIED <id> <saldo>}; \texttt{DISPENSED <id>} $\rightarrow$ \texttt{ACK <id>}; \texttt{BYE}. Regole: il
          bancomat parla per primo, ogni richiesta ha una sola risposta, nessuna operazione prima di un \texttt{OK}, \texttt{<id>} è un
          identificatore unico di transazione scelto dal bancomat.
    \item Bancomat $\rightarrow$ \texttt{AUTH}, banca $\rightarrow$ \texttt{OK}; \texttt{WITHDRAW 17 100} $\rightarrow$ \texttt{GRANTED 17
          420}; il bancomat eroga e invia \texttt{DISPENSED 17} $\rightarrow$ \texttt{ACK 17}; \texttt{BYE}.
    \item Assume un trasporto affidabile e ordinato (TCP), e cifrato e autenticato (TLS). Se la risposta si perde, il bancomat non sa se il
          prelievo è stato autorizzato: con l'identificatore di transazione può ripetere \texttt{WITHDRAW 17} senza che la banca addebiti
          due volte (la banca riconosce il duplicato), e il bancomat eroga solo dopo un \texttt{GRANTED}. Così il cliente riceve i soldi
          \textbf{al più una volta}; se il bancomat si arrende, la banca annulla l'addebito non seguito da \texttt{DISPENSED}.
\end{enumerate}
\end{esercizio}

\begin{esercizio}{FTP, mentre viaggia}
Si avvia una cattura e si apre una sessione anonima su un server FTP pubblico (\texttt{ftp ftp.gnu.org}, utente \texttt{anonymous}),
con \texttt{ls}, \texttt{get README} e \texttt{quit}; filtro \texttt{ftp || ftp-data}.
\begin{enumerate}[label=(\alph*)]
    \item Quante connessioni TCP sono state aperte, verso quali porte del server, e quale è rimasta aperta per tutta la sessione?
    \item Si leggono nome utente e password? Che cosa se ne deduce?
    \item Quale comando e quale risposta hanno preparato ciascuna connessione dati, e chi si è collegato a chi? Come sa il client che il
          file è finito?
\end{enumerate}

\svolgimento
\begin{enumerate}[label=(\alph*)]
    \item Tre: la connessione di \textbf{controllo} verso la porta 21, aperta per tutta la sessione, e una connessione \textbf{dati} per
          \texttt{ls} e una per \texttt{get}, verso porte alte del server.
    \item Sì: \texttt{USER anonymous} e \texttt{PASS} viaggiano in chiaro sulla connessione di controllo. FTP non protegge le credenziali:
          per quelle vere servono FTPS o SFTP.
    \item In modalità passiva il client invia \texttt{PASV} (o \texttt{EPSV}); il server risponde \texttt{227} (o \texttt{229}) con
          indirizzo e porta, ed è il \textbf{client} a collegarsi al server. Il file è finito quando il server chiude la connessione dati
          (FIN): FTP delimita i dati chiudendo la connessione, e la risposta \texttt{226} sulla connessione di controllo lo conferma.
\end{enumerate}
\end{esercizio}

\begin{esercizio}{Il primo chunk di Alice}
Una nuova peer, Alice, entra in un torrent BitTorrent senza alcun chunk. Non avendo nulla da caricare, non può diventare una delle
prime quattro uploader di nessuno. Come ottiene il suo primo chunk?

\svolgimento
Con l'\textbf{optimistic unchoking}: ogni 30 secondi ciascun peer sceglie a caso un vicino in più a cui inviare, indipendentemente da
quanto ne riceve. Prima o poi qualcuno sceglie Alice, che riceve i primi chunk e da lì può cominciare a scambiare.
\end{esercizio}

\begin{esercizio}{Distribuzione client-server e P2P}
Un file di $F = 20$ Gbit deve raggiungere $N$ peer. Il server carica a $u_s = 30$ Mb/s; ogni peer scarica a $d = 2$ Mb/s e carica a $u$.
Con i tempi minimi di distribuzione visti nel corso, riempire una tabella per $N = 10, 100, 1000$ e $u = 300$ kb/s, 700 kb/s, 2 Mb/s, per
client-server e per P2P. Che cosa fa ciascuna architettura al crescere di $N$?

\svolgimento
$D_{cs} = \max\{NF/u_s,\ F/d\}$ e $D_{P2P} = \max\{F/u_s,\ F/d,\ NF/(u_s + Nu)\}$, con $F/u_s \approx 667$ s e $F/d = 10\,000$ s.
\begin{center}
\begin{tabular}{rrrrr}
\toprule
$N$ & client-server & P2P, $u = 300$ kb/s & P2P, $u = 700$ kb/s & P2P, $u = 2$ Mb/s \\
\midrule
10 & 10\,000 s & 10\,000 s & 10\,000 s & 10\,000 s \\
100 & 66\,667 s & 33\,333 s & 20\,000 s & 10\,000 s \\
1000 & 666\,667 s & 60\,606 s & 27\,397 s & 10\,000 s \\
\bottomrule
\end{tabular}
\end{center}
Il client-server cresce \textbf{linearmente} con $N$; il P2P resta limitato, perché ogni peer porta anche capacità di upload: per $N$
grande tende a $F/u$ (e con $u = d$ il collo di bottiglia resta lo scaricamento dei peer).
\end{esercizio}

% ══════════════════════════════════════════════════════════════
\section{HTTP}
% ══════════════════════════════════════════════════════════════

\begin{esercizio}{Vero o falso}
Giustificare ciascuna risposta in una riga.
\begin{enumerate}[label=(\alph*)]
    \item Una pagina fatta di un testo e tre immagini richiede un messaggio di richiesta e quattro messaggi di risposta.
    \item \texttt{www.mit.edu/research.html} e \texttt{www.mit.edu/students.html} possono essere inviate sulla stessa connessione
          persistente.
    \item Con connessioni non persistenti, un solo segmento TCP può portare due messaggi di richiesta HTTP distinti.
    \item L'header \texttt{Date:} di una risposta dice quando l'oggetto è stato modificato l'ultima volta.
    \item Le risposte HTTP non hanno mai il corpo vuoto.
\end{enumerate}

\svolgimento
\begin{enumerate}[label=(\alph*)]
    \item Falso: ogni oggetto ha la sua richiesta, quindi quattro richieste e quattro risposte.
    \item Vero: stanno sullo stesso server, e una connessione persistente porta più oggetti dello stesso server.
    \item Falso: con connessioni non persistenti ogni connessione porta una sola richiesta.
    \item Falso: \texttt{Date:} è il momento in cui la risposta è stata generata; la modifica è in \texttt{Last-Modified:}.
    \item Falso: per esempio \texttt{304 Not Modified}, \texttt{204 No Content} e le risposte a \texttt{HEAD} non hanno corpo.
\end{enumerate}
\end{esercizio}

\begin{esercizio}{Una richiesta, come è stata catturata}
Questi sono i byte esatti di un GET inviato da un browser (\texttt{<cr><lf>} sono i caratteri di fine riga):
\begin{lstlisting}[language=bash]
GET /cs453/index.html HTTP/1.1<cr><lf>
Host: gaia.cs.umass.edu<cr><lf>
User-Agent: Mozilla/5.0 (Windows;U; Windows NT 5.1; en-US; rv:1.7.2) Gecko/20040804 Netscape/7.2<cr><lf>
Accept: text/xml, application/xml, application/xhtml+xml, text/html;q=0.9, */*;q=0.5<cr><lf>
Accept-Language: en-us,en;q=0.5<cr><lf>
Keep-Alive: 300<cr><lf>
Connection: keep-alive<cr><lf>
<cr><lf>
\end{lstlisting}
Dire dove si trova ciascuna risposta: (a) l'URL del documento; (b) la versione di HTTP; (c) connessione persistente o no; (d)
l'indirizzo IP dell'host del browser; (e) il tipo di browser, e perché un server vuole saperlo.

\svolgimento
(a) \texttt{http://gaia.cs.umass.edu/cs453/index.html}: il percorso nella riga di richiesta, l'host nell'header \texttt{Host:}. (b)
HTTP/1.1, alla fine della riga di richiesta. (c) Persistente: \texttt{Connection: keep-alive}. (d) Non c'è: l'indirizzo sta nell'header
IP, non nel messaggio HTTP. (e) \texttt{User-Agent:}, Netscape 7.2 (motore Gecko): il server può inviare versioni diverse dello stesso
oggetto a browser diversi.
\end{esercizio}

\begin{esercizio}{La risposta}
La risposta del server alla richiesta precedente:
\begin{lstlisting}[language=bash]
HTTP/1.1 200 OK<cr><lf>
Date: Tue, 07 Mar 2008 12:39:45 GMT<cr><lf>
Server: Apache/2.0.52 (Fedora)<cr><lf>
Last-Modified: Sat, 10 Dec 2005 18:27:46 GMT<cr><lf>
ETag: "526c3-f22-a88a4c80"<cr><lf>
Content-Length: 3874<cr><lf>
Keep-Alive: timeout=max=100<cr><lf>
Connection: Keep-Alive<cr><lf>
Content-Type: text/html; charset=ISO-8859-1<cr><lf>
<cr><lf>
<!doctype html public "-//w3c//dtd html 4.0 transitional//en"><lf>
\end{lstlisting}
(a) Il documento è stato trovato, e quando è stata generata la risposta? (b) Quando è stato modificato l'ultima volta? (c) Quanti byte
contiene? (d) Quali sono i suoi primi cinque byte? Il server ha accettato una connessione persistente?

\svolgimento
(a) Sì, \texttt{200 OK}; generata il 7 marzo 2008 alle 12:39:45 GMT (\texttt{Date:}). (b) 10 dicembre 2005, 18:27:46 GMT. (c) 3874 byte
(\texttt{Content-Length:}). (d) \texttt{<!doc}. Sì: \texttt{Connection: Keep-Alive}.
\end{esercizio}

\begin{esercizio}{Contare i round trip}
Si clicca su un link. L'indirizzo del server non è in cache: prima che arrivi si visitano $n$ server DNS, con round trip
$\text{RTT}_1, \dots, \text{RTT}_n$. Sia $\text{RTT}_0$ il round trip tra l'host e il server Web; i tempi di trasmissione sono trascurabili.
\begin{enumerate}[label=(\alph*)]
    \item La pagina è un solo piccolo oggetto HTML. Quanto passa dal clic alla ricezione?
    \item L'HTML fa riferimento a otto oggetti piccolissimi sullo stesso server. Quanto ci vuole con HTTP non persistente e nessuna
          connessione parallela?
    \item Con HTTP non persistente e fino a sei connessioni in parallelo?
    \item Con HTTP persistente e pipelining?
\end{enumerate}

\svolgimento
Sia $\Sigma = \text{RTT}_1 + \dots + \text{RTT}_n$.
\begin{enumerate}[label=(\alph*)]
    \item $\Sigma + 2\,\text{RTT}_0$ (un RTT per la connessione TCP, uno per richiesta e risposta).
    \item $\Sigma + 2\,\text{RTT}_0 + 8 \cdot 2\,\text{RTT}_0 = \Sigma + 18\,\text{RTT}_0$.
    \item Otto oggetti in due ondate (sei, poi due): $\Sigma + 2\,\text{RTT}_0 + 2 \cdot 2\,\text{RTT}_0 = \Sigma + 6\,\text{RTT}_0$.
    \item Una sola connessione e tutte le richieste insieme: $\Sigma + 2\,\text{RTT}_0 + \text{RTT}_0 = \Sigma + 3\,\text{RTT}_0$.
\end{enumerate}
\end{esercizio}

\begin{esercizio}{Quanto vale una cache}
Un'istituzione ha una LAN da 100 Mb/s collegata a Internet da un collegamento d'accesso da 15 Mb/s. Gli oggetti sono in media di
1\,000\,000 bit e i browser ne chiedono 16 al secondo. Dal momento in cui il router lato Internet inoltra una richiesta, la risposta
torna in media dopo 3 s. Il tempo di risposta è il ritardo d'accesso più il ritardo di Internet; per l'accesso si usa
$\Delta/(1 - \Delta\beta)$, con $\Delta$ il tempo per inviare un oggetto sul collegamento d'accesso e $\beta$ il tasso di arrivo degli
oggetti.
\begin{enumerate}[label=(\alph*)]
    \item Trovare il tempo di risposta medio totale. Che cosa significa il numero?
    \item Si installa una cache nella LAN con tasso di miss 0,4. Trovare di nuovo il tempo di risposta.
\end{enumerate}

\svolgimento
\begin{enumerate}[label=(\alph*)]
    \item $\Delta = 10^6/(15 \cdot 10^6) \approx 0{,}0667$ s e $\Delta\beta = 0{,}0667 \cdot 16 \approx 1{,}07 > 1$: il collegamento riceve
          più di quanto può inviare, la coda cresce senza limite e il ritardo medio è \textbf{infinito} (la formula non ha senso, ed è
          proprio questo il messaggio).
    \item Al collegamento arrivano solo i miss: $\beta = 0{,}4 \cdot 16 = 6{,}4$, $\Delta\beta \approx 0{,}43$, ritardo d'accesso
          $0{,}0667/(1 - 0{,}43) \approx 0{,}12$ s. Un miss costa $3 + 0{,}12 \approx 3{,}12$ s, un hit circa 0 (10 ms sulla LAN). In media
          $0{,}4 \cdot 3{,}12 + 0{,}6 \cdot 0{,}01 \approx 1{,}25$ s.
\end{enumerate}
\end{esercizio}

\begin{esercizio}{HTTP/2: frame e priorità}
Una pagina è fatta di un video e cinque immagini. Su HTTP/2 il video viaggia in 2000 frame e ogni immagine in 3; il tempo si conta in
tempi di frame.
\begin{enumerate}[label=(\alph*)]
    \item Se si inviano prima tutti i frame del video, senza intercalarli, quanti tempi di frame passano prima che tutte e cinque le
          immagini siano state inviate?
    \item Se i frame sono intercalati, uno per oggetto a turno, quanti?
    \item Ora le immagini hanno priorità sul video, e la prima sulla seconda, e così via. Quanti tempi di frame fino all'invio completo
          della seconda immagine?
\end{enumerate}

\svolgimento
(a) $2000 + 5 \cdot 3 = 2015$. (b) Un giro manda un frame di ciascuno dei sei oggetti; le immagini finiscono al terzo giro: $3 \cdot 6 =
18$. (c) $3 + 3 = 6$.
\end{esercizio}

\begin{esercizio}{Un GET condizionale, scritto a mano}
\texttt{printf} ripete il formato per ogni argomento, quindi ogni riga parte con il \verb|\r\n| richiesto dal protocollo; l'argomento
vuoto è la riga vuota che chiude l'header:
\begin{lstlisting}[language=bash]
$ printf '%s\r\n' 'GET /hypertext/WWW/TheProject.html HTTP/1.1' \
                  'Host: info.cern.ch' 'Connection: close' '' \
  | nc info.cern.ch 80 | head -20
\end{lstlisting}
\begin{enumerate}[label=(\alph*)]
    \item Trovare la riga \texttt{Last-Modified:}. Inviare una seconda richiesta che aggiunga \texttt{If-Modified-Since:} con quella data,
          e ottenere un \texttt{304 Not Modified}. Che cosa c'è nel corpo questa volta?
    \item Quale delle due risposte preferirebbe una cache, e perché?
    \item Rimandare la richiesta condizionale con una data più vecchia, per esempio \texttt{Thu, 01 Jan 1970 00:00:00 GMT}. Perché questa
          volta la risposta è \texttt{200}, e che cosa dice su come decide il server?
\end{enumerate}

\svolgimento
(a) Nulla: un \texttt{304} non ha corpo, solo l'header. (b) Il \texttt{304}: conferma che la copia in cache è ancora valida senza
trasferire di nuovo l'oggetto. (c) Il documento è stato modificato dopo il 1970: il server confronta la data della richiesta con la data
di ultima modifica dell'oggetto, e invia l'oggetto intero solo se quest'ultima è più recente.
\end{esercizio}

\begin{esercizio}{Persistente o no, sul filo}
Si cattura (filtro \texttt{tcp.port == 80}) mentre si scarica due volte la stessa pagina con un solo comando, prima normalmente, poi
chiedendo al server di chiudere dopo ogni risposta:
\begin{lstlisting}[language=bash]
$ u=http://info.cern.ch/hypertext/WWW/TheProject.html
$ curl -s -o /dev/null -o /dev/null $u $u
$ curl -s -H 'Connection: close' -o /dev/null -o /dev/null $u $u
\end{lstlisting}
\begin{enumerate}[label=(\alph*)]
    \item Quanti SYN produce ciascun comando? Chi invia il primo FIN nel secondo caso?
    \item Misurare l'RTT dall'handshake. Per ogni oggetto, quanti RTT passano dal primo pacchetto della sua connessione (o dal suo GET) al
          primo byte della risposta?
    \item I numeri tornano con il conteggio fatto nel corso per connessioni persistenti e non persistenti?
\end{enumerate}

\svolgimento
(a) Il primo comando un solo SYN (una connessione persistente per i due oggetti); il secondo due. Con \texttt{Connection: close} chiude
per primo il \textbf{server}, dopo aver inviato la risposta. (b) Persistente: il primo oggetto arriva dopo 2 RTT dal SYN, il secondo
dopo 1 RTT dal suo GET; non persistente: 2 RTT ciascuno. (c) Sì: è il conteggio $2\,\text{RTT}$ per oggetto contro $\text{RTT}$ per
oggetto dopo il primo.
\end{esercizio}

% ══════════════════════════════════════════════════════════════
\section{Posta e DNS}
% ══════════════════════════════════════════════════════════════

\begin{esercizio}{Il percorso di un messaggio}
Alice invia un messaggio da un account di webmail; Bob legge la posta con IMAP. Elencare in ordine i protocolli applicativi che il
messaggio attraversa dall'host di Alice a quello di Bob, e tra quali macchine gira ciascuno.

\svolgimento
HTTP(S) tra il browser di Alice e il suo server di webmail; SMTP tra il server di posta di Alice e quello di Bob; IMAP tra il client di
Bob e il suo server di posta.
\end{esercizio}

\begin{esercizio}{La fine del corpo}
Come segna SMTP la fine del corpo di un messaggio? E HTTP? HTTP potrebbe usare il metodo di SMTP? Perché?

\svolgimento
SMTP con una riga che contiene solo un punto (\verb|\r\n.\r\n|). HTTP con \texttt{Content-Length:} (o con la codifica
\emph{chunked}, o chiudendo la connessione). HTTP non può usare il punto: il suo corpo può essere binario e contenere quella sequenza,
quindi servirebbe uno stuffing di tutti i dati (SMTP lo fa raddoppiando i punti a inizio riga, ma lavora su testo).
\end{esercizio}

\begin{esercizio}{Chi ha mentito?}
Ogni server di posta che tratta un messaggio aggiunge in cima agli header una riga \texttt{Received:}, con chi gli ha consegnato il
messaggio e chi l'ha preso. Questi sono gli header di uno spam, dal più recente:
\begin{lstlisting}[language=bash]
Received: from barmail.cs.umass.edu ([128.119.240.3]) by cs.umass.edu; 7 Nov 2008 13:27:10 -0500
Received: from asusus-4b96 (localhost [127.0.0.1]) by barmail.cs.umass.edu (Spam Firewall); 7 Nov 2008 13:27:07 -0500
Received: from asusus-4b96 ([58.88.21.177]) by barmail.cs.umass.edu; 7 Nov 2008 13:27:07 -0500
Received: from [58.88.21.177] by inbnd55.exchangeddd.com; 8 Nov 2008 01:27:07 +0700
From: "Jonny" <tennis5@pp33head.com>
To: <hg@cs.umass.edu>
\end{lstlisting}
\begin{enumerate}[label=(\alph*)]
    \item Solo chi ha originato il messaggio è malevolo; tutti gli altri host sono onesti. Quale host ha generato il messaggio? Di quali
          righe ci si può fidare?
    \item Guardare gli header di un messaggio ricevuto di recente: quante righe \texttt{Received:} ci sono, e si trova l'indirizzo IP
          dell'host del mittente?
\end{enumerate}

\svolgimento
(a) L'ha generato l'host \texttt{58.88.21.177}. Le prime tre righe sono state scritte da server onesti: la terza in particolare è di
\texttt{barmail}, che ha visto la connessione arrivare da \texttt{58.88.21.177} (il nome \texttt{asusus-4b96} invece lo dichiara il
mittente, e non è affidabile). L'ultima riga, che parla di \texttt{inbnd55.exchangeddd.com}, sta \emph{sotto} il primo server onesto:
l'ha scritta l'originatore stesso, ed è falsa. (b) Di solito da due a sei righe; l'indirizzo dell'host del mittente si trova nella riga
più in basso scritta dal suo primo server, se il provider non lo nasconde.
\end{esercizio}

\begin{esercizio}{Lo stesso alias per Web e posta}
Il server Web e il server di posta di un'organizzazione possono avere esattamente lo stesso alias, per esempio \texttt{foo.com}? Quale
tipo di resource record contiene il nome del server di posta?

\svolgimento
Sì: il server Web si trova con un record \texttt{A} (o \texttt{CNAME}) per \texttt{foo.com}, il server di posta con un record
\texttt{MX} per lo stesso nome, che contiene il nome canonico del server di posta.
\end{esercizio}

\begin{esercizio}{La cache del DNS di dipartimento}
Il dipartimento ha un server DNS locale per tutte le sue macchine, e voi siete un utente normale, non un amministratore. Potete scoprire
se un sito esterno è stato visitato da una macchina del dipartimento pochi secondi fa?

\svolgimento
Sì: si interroga il server locale \textbf{senza ricorsione} (\texttt{dig +norecurse @server nome}). Se risponde con il record, ce l'ha
in cache, quindi qualcuno l'ha chiesto di recente; il TTL rimasto, più basso di quello originale, dice anche quanto tempo fa. Se non
risponde, nessuno l'ha chiesto entro il TTL.
\end{esercizio}

\begin{esercizio}{Giù per l'albero a mano}
Si fa ciò che fa il server locale con la cache vuota: si parte da un root server e si seguono i rinvii fino al server autoritativo per
il nome.
\begin{lstlisting}[language=bash]
$ dig +norecurse @a.root-servers.net www.unina.it A      # poi si chiedono i server che nomina
\end{lstlisting}
\begin{enumerate}[label=(\alph*)]
    \item Scrivere la catena di deleghe: ogni server interrogato, la zona per cui risponde e i server a cui rinvia.
    \item Ripetere per un nome sotto un dominio di primo livello diverso, per esempio un'università giapponese.
    \item Nell'ultima risposta, il flag \texttt{aa} è acceso, e qual è il TTL? Chiedere lo stesso nome al resolver locale: quale TTL è più
          piccolo, e perché?
    \item Qualche passo ha risposto con un CNAME? Dove ha portato la catena?
\end{enumerate}
Prima si spegne qualsiasi VPN: una VPN può rispondere al posto del root server. Un root server accende \texttt{aa} solo se è autoritativo
e non accende mai \texttt{ra}; se si vede \texttt{ra}, qualcuno sulla strada sta rispondendo al suo posto.

\svolgimento
\begin{enumerate}[label=(\alph*)]
    \item Root server (zona \texttt{.}) $\rightarrow$ rinvio ai server di \texttt{it.} (\texttt{*.nic.it}, \dots) $\rightarrow$ i
          server di \texttt{it.} rinviano ai name server di \texttt{unina.it.} $\rightarrow$ questi danno la risposta. Ogni rinvio ha la
          sezione AUTHORITY con i record \texttt{NS} e, se ci sono, i \emph{glue} nella sezione ADDITIONAL.
    \item Stessa forma: root $\rightarrow$ \texttt{jp.} $\rightarrow$ \texttt{ac.jp.} $\rightarrow$ il dominio dell'università: un
          livello in più.
    \item Sì, \texttt{aa} è acceso nella risposta del server autoritativo, con il TTL pieno impostato dalla zona. Il resolver locale
          risponde con un TTL più piccolo (o uguale, se l'ha appena chiesto): lo prende dalla sua cache, che fa scendere il TTL con il passare
          del tempo.
    \item Spesso \texttt{www} è un CNAME verso un nome di una CDN o di un altro dominio: la catena allora riparte dalla zona di quel nome.
\end{enumerate}
\end{esercizio}
